\chapter{The Slip Becomes A Number}

\section{A crossing predicate}

A crossing predicate is a finite question, not yet a root or a physical event.  Let the positive rational domain be
\[
 \mathbb Q_{>0}^{\rm rec}
 =\{(p,q):p,q\in\mathbb N,\ p>0,\ q>0\},
 \qquad d=p/q.
\]
Suppose a report is represented by integers \(A(p,q)\ge0\) and \(B(p,q)>0\), with reported ratio \(A/B\), and let
target \(t\ge0\) be an integer in compatible declared units.  The exact crossing question is
\[
                  C_t(p,q):=[t\,B(p,q)\le A(p,q)].
\]
It uses integer cross-multiplication.  No decimal approximation or floor is required to decide the branch.

A forward evaluator constructs
\[
 E_t(p,q)=(A(p,q),B(p,q),tB(p,q),A(p,q)).
\]
The requirement question on this report asks whether the third entry is at most the fourth.  Reading pulls that question
back through \(E_t\); the evaluator is covariant and the requirement read is contravariant.  The question does not
construct a distance or choose a side of a bracket.

For the instrument's finite crossing model,
\[
                    g(d)=\frac{18}{d^2},\qquad t=5.
\]
At \(d=p/q>0\),
\[
              C_5(p,q):=[5p^2\le18q^2].
\]
The constants \(18\) and \(5\), the positive domain, and the interpretation of \(d\) are posited calibration data.
If units are attached, \(18/d^2\) and \(5\) must have matching units before comparison.

The monotone orientation follows directly.  For \(0<d_1\le d_2\),
\[
 d_1^2\le d_2^2,\qquad
 \frac1{d_1^2}\ge\frac1{d_2^2},\qquad
 \frac{18}{d_1^2}\ge\frac{18}{d_2^2}.
\]
Therefore the yes region for \(C_5\) lies downward in distance.  The floored report
\(\lfloor18/d^2\rfloor\) is also nonincreasing, but has plateaus on which distinct ratios receive the same integer.

A bracket claim needs additional premises.  On an ordered domain \(D\), require:

\[
 a<b,\qquad C_t(a)=\mathsf{yes},\qquad C_t(b)=\mathsf{no},
\]
and monotonicity in the declared orientation,
\[
 d_1<d_2\quad\Longrightarrow\quad g(d_1)\ge g(d_2).
\]
These facts certify that the bracket straddles a truth transition.  They do not automatically prove that an exact
rational equality \(g(d)=t\) exists inside; the boundary may lie outside the represented rational set.

A finite scan makes the method auditable.  Given ordered candidates
\(D_N=[d_1,\ldots,d_N]\) and budget \(B\le N\), visit \(i=1,\ldots,B\).  Verify the domain conditions, construct the
cross-products, resolve \(C_t(d_i)\), and append the full row to crossed, not-crossed, or domain-error slices.  Stop with
\textsc{complete}, \textsc{budget}, \textsc{domain-error}, or \textsc{resolver-domain}.  The certificate contains
candidate indices and rationals, unreduced and reduced forms when used, cross-products, decisions, slices, units,
budget, completeness, erasures, and stop reason.  An unvisited suffix carries no decision.

\paragraph{Worked example 1: exact endpoint crossing (instrument computation).}
At \(d=1=1/1\),
\[
                         5\cdot1^2=5\le18\cdot1^2=18,
\]
so the predicate answers yes.  At \(d=2=2/1\),
\[
                         5\cdot2^2=20>18,
\]
so it answers no.  The ordered endpoints straddle the transition.  At \(d=3/2\),
\[
                         5\cdot3^2=45\le18\cdot2^2=72,
\]
so the midpoint candidate remains on the crossed side.

\paragraph{Worked example 2: a finite rational slice (instrument computation).}
Visit \(d=9/5,19/10,2\).  The exact comparisons are
\[
\begin{array}{c|c|c}
d&5p^2&18q^2\\
\hline
9/5&405&450\\
19/10&1805&1800\\
2/1&20&18
\end{array}
\]
and the decision slice is
\[
 \text{crossed}=[9/5],\qquad
 \text{not-crossed}=[19/10,2].
\]
Thus the sampled transition lies between \(9/5\) and \(19/10\).  This is a statement about the sampled exact
predicate, not yet a certified narrowest bracket.

A scan certifies only the visited rows and, when complete, the transitions between adjacent listed candidates.  It does
not produce the canonical seed-free continued-fraction jar.  The canonical distance walls begin with \(17/9\) and
\(2\):
\[
 17/9<2,\qquad
 5\cdot17^2=1445\le1458=18\cdot9^2,\qquad
 5\cdot2^2=20>18.
\]
Thus \(17/9\) is on the yes side and \(2\) on the no side.  Their later continued-fraction role may be developed in
the jar section; the exact opposed-wall arithmetic already belongs in this certificate.

Monotonicity cannot be inferred from one change of branch.  It must be established from the report formula or tested
only as a finite empirical pattern.

\paragraph{Worked example 3: nonmonotone counterexample (external numerical experiment).}
Let
\[
                         h(d)=(d-1)(d-2),
 \qquad C(d)=[h(d)\ge0].
\]
The predicate is yes at \(d=1/2\), no at \(d=3/2\), and yes again at \(d=3\).  A scan sees two transitions.  Keeping one
half after a branch test without a monotonicity or sign-change argument could discard another crossing.  The shared word
crossing does not license a bisection invariant.

\paragraph{Worked example 4: malformed rational record (computation).}
The record \((p,q)=(3,0)\) has no represented positive rational value.  Cross-multiplication might still produce
integers, but the domain certificate fails first.  The correct result is \textsc{domain-error}, not a crossing decision.
Likewise, a negative or unit-incompatible target requires a different declared domain and question.

A compact rational bisection contract is also available.  Validate positive ordered endpoints
\[
 \ell<u,\qquad C_t(\ell)=\mathsf{yes},\qquad C_t(u)=\mathsf{no}.
\]
Form the exact rational midpoint \(m=(\ell+u)/2\).  If \(C_t(m)\) is yes, replace \(\ell\) by \(m\); otherwise replace
\(u\) by \(m\).  Each accepted step preserves endpoint order and the opposed-bit invariant.  With finite fuel \(N\),
the certificate returns the final rational bracket, every midpoint and branch decision, unused fuel, and stop reason.
It certifies a truth-transition bracket, not an exact represented root.  This legacy fuel-bisection route is distinct
from the canonical seed-free continued-fraction jar; agreement between their walls requires a separate comparison.

The floor report is related but distinct.  For nonnegative \(A\) and positive \(B\),
\[
                         \left\lfloor\frac AB\right\rfloor\ge t
 \quad\Longleftrightarrow\quad
                         tB\le A.
\]
This identity explains why the exact crossing question can be evaluated without performing division.  The certificate
must retain \(A,B,tB\), and the comparison result; storing only the Boolean branch erases the distance from threshold.
Equal floors do not imply equal ratios or equal distances.  The fractional residue
\[
                         A-B\left\lfloor A/B\right\rfloor
\]
records information erased by the floor.  Also, unreduced pairs such as \((p,q)\) and \((2p,2q)\) represent the same
rational distance while remaining different records; canonicalization by a positive greatest common divisor must be
declared when record equality is intended to mean rational equality.

\paragraph{Worked example 5: static and kinetic slip (physics, then interpretation).}
A friction experiment may define a physical breakaway event through applied force, normal load, surface preparation,
velocity, temperature, and an uncertainty model.  Physics owns those variables and the distinction between static and
kinetic regimes.  Reading \(C_t(d)\) as that event is an interpretation requiring an apparatus map from rational
distance records to measured forces.  The finite inequality alone supplies no friction law, discontinuity, or
unpredictability claim.

The crossing audit is
\[
 \mathcal X=(D,t,\text{units},A,B,C_t,\text{monotonicity premise},
             \text{endpoint orientation/bits},\text{raw }A/B,\text{floor residue},
             \text{canonicalization policy},D_N,N,\text{budget/fuel},
             \text{cross-products},\text{decisions},\text{slices},
             \text{continued-fraction or bisection route},\text{unresolved frontier},
             \text{integer arithmetic assumption},\text{completeness},\text{stop},
             \text{erasures},\text{provenance}).
\]
An absent monotonicity argument, failed domain check, incompatible unit, partial scan, or unresolved arithmetic
operation is recorded directly.  The audit must state whether integers are unbounded; for bounded integers it must
instead certify that every square and cross-product avoids overflow.

The action ledger closes the section.  Domain, target, units, report formula, candidate order, canonicalization,
arithmetic model, route, budgets, and stop policies are \emph{posited}.  Monotonicity, cross-product identities,
endpoint inequalities, bisection invariants, finite slice coverage, and implications proved from declared premises are
\emph{demonstrated}.  Actual candidate reports, floors, residues, midpoint and branch decisions, slices, budgets spent,
frontiers, and stops are \emph{read}.  Existence of an exact represented root, uniqueness,
physical breakaway, and claims beyond the scanned prefix remain \emph{reserved}.  A crossing predicate supplies one
finite question; bracketing requires the premises that let its answers be sliced safely.

\section{Bisection is iteration by slicing}

A bisection bracket is finite data:
\[
 \mathcal B=(\ell,u,c_\ell,c_u),
 \qquad \ell,u\in\mathbb Q_{>0}^{\rm rec},
\]
where \(c_\ell=C_t(\ell)\) and \(c_u=C_t(u)\).  In the downward-yes orientation, a valid input satisfies
\[
                         \ell<u,\qquad c_\ell=\mathsf{yes},
                         \qquad c_u=\mathsf{no}.
\]
The bracket does not contain a point value or a guarantee that a represented equality exists.

For \(\ell=p/q\) and \(u=r/s\), with positive denominators, the exact midpoint is
\[
                         m=\frac{ps+rq}{2qs}.
\]
A one-step slice evaluates \(c_m=C_t(m)\) and returns
\[
 \operatorname{bisectOnce}(\mathcal B)=
 \begin{cases}
 ((\ell,u,m,c_m),(m,u,\mathsf{yes},\mathsf{no})),&c_m=\mathsf{yes},\\
 ((\ell,u,m,c_m),(\ell,m,\mathsf{yes},\mathsf{no})),&c_m=\mathsf{no}.
 \end{cases}
\]
The first component is the step ledger; the second is the retained half.  Exact arithmetic may leave the midpoint
unreduced, so a canonicalization policy must be declared separately.

The pure loop is fuel-only and recursive:
\[
\begin{aligned}
 \operatorname{run}(0,\mathcal B)&=(\mathcal B,[]),\\
 \operatorname{run}(n+1,\mathcal B)&=
 \text{let }(\sigma,\mathcal B')=\operatorname{bisectOnce}(\mathcal B),\\
 &\quad\text{let }(\mathcal B'',T)=\operatorname{run}(n,\mathcal B')
 \text{ in }(\mathcal B'',\sigma::T).
\end{aligned}
\]
Thus zero fuel returns the current bracket and an empty trace.  Successor fuel performs one step, recurses on the
updated bracket with the remaining fuel, and prepends the current step to the remaining trace, preserving chronological
order.  A successful run with fuel \(N\) has trace length \(N\).  Width-target stopping is a proposed extension, not
part of this pure recurrence.

The certificate contains initial and final brackets, every midpoint, raw cross-products, branch decisions, canonical
forms when used, fuel and trace length, widths, erasures, and stop reason.  A fuel-complete result owns a finite bracket,
not an unread infinite tail.

A replay checker validates each adjacent row: recompute the exact midpoint; recompute cross-products and crossing bit;
verify that the selected branch produces the recorded next bracket; verify that the next row's old bracket equals that
result; recheck opposed endpoint bits, order, and width; and finally check fuel accounting and trace length.  Replay
success certifies adjacency and invariant preservation for the recorded finite run.

The familiar width law is conditional.  Interpret the positive rationals in an ordered metric line, use the exact
arithmetic midpoint, preserve endpoint order and opposed crossing bits, and assume the monotone enclosure orientation
from the previous section.  Then
\[
       w_k=u_k-\ell_k,\qquad
       w_{k+1}=\frac{w_k}{2},\qquad
       w_n=w_0\,2^{-n}.
\]
This follows by induction from retaining one of two equal halves.  The formula proves shrinking width under those
hypotheses.  It does not prove an exact root, uniqueness without stricter premises, or physical correctness.

\paragraph{Worked example 1: four exact slices (instrument computation).}
For \(C_5(p,q)=[5p^2\le18q^2]\), begin with \([1,2]\).  The exact trace is
\[
\begin{array}{c|c|c|c}
k&[\ell_k,u_k]&m_k&C_5(m_k)\\
\hline
0&[1,2]&3/2&\mathsf{yes}\\
1&[3/2,2]&7/4&\mathsf{yes}\\
2&[7/4,2]&15/8&\mathsf{yes}\\
3&[15/8,2]&31/16&\mathsf{no}
\end{array}
\]
and the retained bracket is
\[
                         [15/8,31/16].
\]
Its width is \(1/16=(2-1)2^{-4}\).  The certificate owns four cross-multiplication decisions and this bracket; it does
not own the irrational equality boundary \(\sqrt{18/5}\) as a represented rational point.

\paragraph{Worked example 2: malformed enclosure (computation).}
Take \([1,3/2]\).  Both endpoints answer yes.  A midpoint can still be computed, but there are no opposed endpoint bits,
so the bisection contract stops with \textsc{invariant-failure}.  Repeated halving of an arbitrary interval is not a
crossing enclosure.

\paragraph{Worked example 3: opposed bits without monotonicity (external counterexample).}
Let
\[
                     h(d)=-(d-1)(d-2)(d-3),\qquad C(d)=[h(d)\ge0].
\]
On \([1/2,7/2]\), the lower endpoint is yes and the upper endpoint no, but three crossings lie inside.  The predicate is not
monotone.  Binary slicing may retain a half containing one transition, yet it cannot claim a unique boundary or safely
apply the downward-yes invariant globally.  Opposed bits alone are insufficient.

Endpoint projection is lossy.  Define \(\pi_{\rm final}\) to retain only the final \((\ell_N,u_N)\).  Distinct step
histories can share that pair while differing in raw midpoint records, reductions, costs, or repeated checks.  A
question on final endpoints cannot recover those histories.

\paragraph{Worked example 4: same bracket, different ledger (computation).}
One run may reduce every midpoint immediately; another may retain unreduced pairs and canonicalize only the final
endpoints.  If both branch decisions agree, their rational final brackets may be equal while their stored traces and
costs differ.  The audit must retain or explicitly erase canonicalization history; endpoint equality does not identify
the algorithms.

Newton's method is a different route:
\[
                         x_{k+1}=x_k-\frac{f(x_k)}{f'(x_k)}.
\]
It requires a differentiable model, nonzero derivative at each used point, a starting value, and its own stopping
argument.  It need not preserve an enclosure.  Quasi-Newton methods replace exact derivative information with updated
approximations and require still different premises.  Their speed or endpoint agreement does not transfer bisection's
opposed-bit invariant.

\paragraph{Worked example 5: Newton step leaves an interval (external counterexample).}
For \(f(x)=x^3-2x+2\), start at \(x_0=0\).  Then
\[
                         x_1=0-\frac{2}{-2}=1,\qquad
                         x_2=1-\frac{1}{1}=0,
\]
so the route cycles between \(0\) and \(1\).  This does not make Newton invalid; it shows that tangent iteration and
bisection slicing have different guarantees and failure certificates.

The canonical jar route is different again.  It uses the seed-free continued-fraction algorithm: Euclidean
division-with-remainder produces whole partial quotients, and their convergents produce the ordered distance walls
\(17/9\) and \(2/1\).  It is not midpoint bisection, a raw mediant crawl, or a grid search.  Those walls map to the
canonical count-three report jar written \([129.6,137.7]\).  The rational distance bisection above is a legacy
fuel-driven route.  Neither route inherits the other's certificate; wall agreement needs an explicit comparison map.

The bisection audit is
\[
 \mathcal B_{\rm audit}=(C_t,\text{domain},\text{orientation},\mathcal B_0,
 \text{monotonicity/enclosure evidence},N,\text{proposed width-target extension},
 \text{midpoints},\text{cross-products},\text{decisions},
 \text{canonicalization},\text{bracket trace},\text{widths},
 \text{fuel},\text{trace length},\text{replay checks},\text{frontier},\text{stop},\text{erasures},
 \text{route identity},\text{provenance}).
\]

The action ledger is exact.  Domain, crossing question, initial bracket, orientation, arithmetic model,
canonicalization, fuel, and any proposed width-target extension are \emph{posited}.  Midpoint identities, recursive
trace-length law, replay-checker implications, invariant preservation, conditional width law, and final-bracket
implications are \emph{demonstrated}.  Actual midpoints, cross-products,
branch decisions, bracket trace, fuel use, widths, and stop are \emph{read}.  Exact root existence, uniqueness,
Newton or quasi-Newton guarantees, continued-fraction wall identity, and physical slip meaning remain
\emph{reserved}.  Bisection is iteration by slicing because every update keeps one evidenced half and records the one
it discarded.

\section{A report is not a certificate}

A numerical routine may return data without establishing every claim a reader wants to attach to it.  Separate four
layers:
\[
\begin{array}{rcl}
\textit{return}&::=&\operatorname{none}(\textit{reason})
 \mid\operatorname{some}(\textit{payload}),\\
\textit{report}&::=&\operatorname{report}(\textit{schema},\textit{route},
 \textit{inputs},\textit{outputs},\textit{trace},\textit{stop}),\\
\textit{validation}&::=&\operatorname{accepted}\mid
 \operatorname{rejected}(\textit{failures})\mid\operatorname{unresolved},\\
\textit{certificate}&::=&\operatorname{certificate}(
 \textit{report},\textit{checks},\textit{claim},\textit{provenance}).
\end{array}
\]
An optional return distinguishes absence from a numerical zero or empty bracket.  A report organizes returned data.
Validation checks a declared schema and claim.  Only an accepted validation can be packaged as a certificate for that
specific claim.

The conversion chain is explicitly one-way:
\[
 \text{raw bytes or text}\longrightarrow\text{parsed report}
 \longrightarrow\text{validated report}\longrightarrow\text{certified claim}.
\]
Raw input does not imply successful parsing; parsing does not imply validity; validity does not imply every possible
claim; and certification of one claim does not certify another.

If a Boolean validator \(V\) is used for property \(P\), acceptance licenses \(P(r)\) only when the soundness
implication
\[
                         V(r)=\mathsf{true}\Longrightarrow P(r)
\]
has been demonstrated or is explicitly trusted or posited.  The actual Boolean outcome is \emph{read}.  A validator
name or true bit cannot supply its own soundness argument.

The reporting pipeline is finite:

\[
\begin{array}{ll}
1.&\text{acquire the optional return within budget},\\
2.&\text{parse schema and version},\\
3.&\text{validate domains, units, scales, and route identity},\\
4.&\text{replay transitions and arithmetic},\\
5.&\text{evaluate the requested claim and slice checks into pass/fail/unresolved},\\
6.&\text{emit a certificate only when every required check passes.}
\end{array}
\]

Stops are \textsc{none}, \textsc{budget}, \textsc{parse-error}, \textsc{version-mismatch},
\textsc{schema-unknown}, \textsc{required-entry-missing}, \textsc{unit-scale-error},
\textsc{domain-error}, \textsc{arithmetic-resource}, \textsc{overflow},
\textsc{partial-report}, \textsc{route-mismatch}, \textsc{replay-failure},
\textsc{certificate-missing}, \textsc{certificate-invalid}, \textsc{claim-failure}, and \textsc{complete}.  The receipt records raw
return, parsing decisions, all validation rows, budget spent, checked scope, unresolved frontier, erasures, stop reason, and provenance.
A failed or partial validation does not erase the report; it limits what the report licenses.

For a bracket report, validation questions include
\[
 \ell<u,\qquad C_t(\ell)=\mathsf{yes},\qquad C_t(u)=\mathsf{no},
\]
together with domain and arithmetic assumptions.  If a trace is supplied, replay checks every midpoint, branch,
adjacent bracket, fuel count, and final endpoint.  If no trace is supplied, endpoint checks may still pass, but no claim
about how the endpoints were reached is certified.

\paragraph{Worked example 1: a replayable bisection report (computation).}
A four-step report contains
\[
 [1,2]\to[3/2,2]\to[7/4,2]\to[15/8,2]\to[15/8,31/16].
\]
The validator recomputes midpoints \(3/2,7/4,15/8,31/16\), checks their crossing bits, verifies every retained half,
and confirms trace length four and final width \(1/16\).  If all rows pass, the certificate owns this finite bracket and
trace.  It does not own an exact represented crossing point.

\paragraph{Worked example 2: a malformed bracket report (computation).}
Suppose a payload reports
\[
                         \ell=2,\qquad u=1.
\]
Both values parse as positive rationals, but the order check fails.  The result is a valid parsed report and a rejected
bracket certificate.  Swapping the endpoints during validation would silently change the returned data; a repair must
be a new, explicitly transformed report with its own provenance.

\paragraph{Worked example 3: absence is not a default value (computation).}
If acquisition returns \(\operatorname{none}(\textsc{budget})\), the parser must not synthesize
\([0,0]\), a zero endpoint, or an empty trace and call it a reading.  The receipt owns only the budget stop and unresolved
claim.  Such a none result is scoped to the scanned candidates, budget, and run configuration; it is not a global
nonexistence statement.  A later run may produce data, but it is a separate acquisition.

A summary projection also loses evidence.  Let
\[
 \pi_{\rm endpoints}(\text{full report})=(\ell,u).
\]
Distinct traces may share endpoints.  Questions about endpoint order can be pulled through this projection, but
questions about midpoint history, cost, or route cannot be answered after those entries are erased.  Equal summaries do
not identify full reports.

\paragraph{Worked example 4: jar ownership versus endpoint diagnostics (instrument computation).}
The canonical count-three report owns the ordered-wall jar
\[
                         [129.6,137.7].
\]
Two finer route diagnostics report, at scale \(10^{18}\),
\[
 d=137011290548979455469,\qquad
 s=137011290753751157155,
\]
with scaled difference \(204771701686\).  Those diagnostics support a declared-resolution comparison, but neither
replaces the jar.  A valid report tags the continued-fraction jar route, direct route, and stepped route separately;
shared digits do not merge their traces or costs.

The integers are unequal, yet coarse formatting maps both to the displayed string \(137.011\).  This is a display
collision, not equality of readings.  A loss-auditable decimal or scaled-floor report retains the raw integer or exact
rational, scale, formatting rule, displayed quotient, and discarded residue.  For exact ratio \(A/B\) at scale \(M\),
\[
 n=\left\lfloor\frac{AM}{B}\right\rfloor,\qquad
 \rho=AM-nB.
\]
The floor erases \(\rho/B\); retaining \(A,B,M,n,\rho\) permits replay.

\paragraph{Worked example 4b: flags are not a certificate (computation).}
A payload containing entries \(\mathtt{verified=true}\) and \(\mathtt{complete=true}\) may parse under a permissive
schema.  If the required certificate is absent, validation returns \textsc{certificate-missing}; if its checks fail, it
returns \textsc{certificate-invalid}.  Self-reported flags are data, not validation evidence.

\paragraph{Worked example 5: same payload, different route (computation).}
Suppose two routines return the rational pair \((3,2)\).  One obtained it by reducing \(18/12\); the other read it
directly from an input table.  The endpoint payloads agree, but the first certificate may own a Euclidean remainder
trace while the second owns a table lookup.  A claim about value equality can pass; a claim about equal process or cost
must fail or remain unresolved.

Physical reporting requires another provenance chain.  A displayed voltage interval is not automatically a calibrated
measurement interval.  Physics owns sensor response, units, calibration transfer, environmental conditions, sampling
plan, and uncertainty model.  Numerical validation can certify formatting, arithmetic, and finite acquisition records;
it cannot manufacture empirical coverage or physical interpretation.

A reproducibility receipt records the checker and compiler releases, dependency lock and versions, source-snapshot
hash, run and build configuration, arithmetic model, platform, protocol version, run identifier and time, deterministic
inputs, and every declared nondeterministic input---including seeds, schedules, and environmental state where relevant.
Reproduction is conditional on matching that declared provenance.  Agreement under one configuration establishes
neither physical reproducibility nor invariance across platforms, arithmetic models, compilers, or checker releases.

The report audit is
\[
\begin{aligned}
\mathcal R={}&(\text{raw bytes/text and optional return},\text{encoding},\text{schema/version},\text{parser},\\
&\text{validator/checker identity and release},\text{soundness status},\text{compiler release},
\text{dependency lock/versions},\\
&\text{source-snapshot hash},\text{run/build configuration},\text{arithmetic model},\text{platform},
\text{protocol},\text{run ID/time},\\
&\text{deterministic and declared nondeterministic inputs},\text{route identity},\text{units/scale},
\text{exact raw values},\\
&\text{formatting rule, quotient, and residue},\text{outputs},\text{trace},\text{budget},
\text{checked scope},\text{unresolved frontier},\text{stop},\\
&\text{domain/arithmetic/replay checks},\text{failures},\text{claim syntax/evaluator},
\text{pass/fail/unresolved slices},\\
&\text{summary projections},\text{erasures},\text{certificate presence and validity},
\text{certified claims},\text{provenance}).
\end{aligned}
\]

\paragraph{Descent experiment: a boundary condition collapses two branches to one.}
The Dirichlet--Bancroft experiment treats a crossing report as a selection over an explicit finite candidate set rather than as a promise of uniqueness.
Its candidate domain is \(X=\{\mathsf{nearEarth},\mathsf{deepSpace}\}\), its constraint domain is
\[
 C=\{\mathsf{geometryOnly},\mathsf{earthBoundBoundary}\},
\]
and its transform returns a ledger,
\[
 T:X\times C\longrightarrow\mathsf{Ledger}(X),
 \qquad
 T(X,c)=\operatorname{filter}(\operatorname{admits}(c),X).
\]
Equivalently, the executable setup consumes only \(c\) because the two-element candidate list is fixed and visible in the experiment.
Under geometry alone both branches pass and the run returns count \(2\); under the earth-bound boundary only the near-Earth branch passes and the run returns \(1\).
The exported claim is the boundedness predicate \(\#T(X,c)\le1\), so the canonical boundary setup succeeds and the geometry-only counterexample fails.
The trace is two Boolean constraint applications, one finite filter, and one length computation; rejected branches
remain recorded.
Its precondition is that the candidate list is the declared search universe.
Exhausting those two candidates proves nothing about an omitted third solution, and uniqueness within the list is therefore the stop rule's exact ceiling.
The commuting test compares filtering then counting with the explicit admissible-branch enumeration; the residual test is the branch removed by the stronger boundary, which remains named in the rejected slice.
Satellite clocks, range equations, and a receiver are absent; the names denote boundary-conditioned finite selection,
not the continuous Dirichlet or Bancroft equations.

\paragraph{Descent experiment: the native apparatus brackets a sign crossing.}
The positron-threshold experiment is categorically different because it imports \texttt{Measurement.Apparatus} and consumes the apparatus's own paths, variation, second-difference operator, and event ledgers.
Let \(P=\mathsf{CubicGaugePath0}\), \(V=\mathsf{CubicGaugeVariation0}\), and define the native transform
\[
 \Delta^2_{\rm pair}:P\times V\longrightarrow\mathbb Z,
 \qquad
 \Delta^2_{\rm pair}(p,v)=\operatorname{pairDelta2}(p,v).
\]
The consumed variation is \(v=\operatorname{pairVariation}(\mathit{node1},\mathit{node2})\).
The two consumed paths are the native flat path at strain tag \(0\) and tilted path at tag \(1\), and the apparatus computes
\[
 d_0=\Delta^2_{\rm pair}(\mathit{flatPath},v)=-1,
 \qquad
 d_1=\Delta^2_{\rm pair}(\mathit{tiltedPath},v)=1.
\]
Independently, the apparatus's event transform supplies \(\operatorname{antimatterCount}(\mathit{flatEvents})=0\) and \(1\le\operatorname{antimatterCount}(\mathit{tiltedEvents})\).
These are two neighboring finite probes, so the only demonstrated admission bracket is \((0,1]\).
The knot transform consumes both signed readings,
\[
 K(s_0,d_0,s_1,d_1)
 =\bigl(s_0(d_1-d_0)-d_0(s_1-s_0),\;d_1-d_0\bigr),
\]
and returns the numerator--denominator pair \(K(0,-1,1,1)=(1,2)\), representing the zero of the secant through the two knots.
The top-level claim and its supporting theorems consume the two native paths, pair variation, second differences, event counts, and exact integer knot arithmetic.
The exported run is narrower: one setup selects either the flat or tilted event list and returns only its \texttt{antimatterCount}; it does not evaluate \(\Delta^2_{\rm pair}\) or \(K\).
The source emits no failure protocol or adaptive search: its ceiling is the two fixed probes and separately proved
knot pair.
The comparison test places the sign of \(\Delta^2_{\rm pair}\) beside the charge-event slice; it does not prove a general equivalence between them.
The residual test retains the full before/after pair, because the rational crossing \(1/2\) is an interpolation and not itself an apparatus event.
These are compiler evaluations of the native apparatus dependency graph, not a physical positron observation,
detector threshold, efficiency, or unit.  The local \texttt{deviceNear} sees only the selected count predicate; it does
not report the distinct second-difference values or knot pair.

The action ledger is exact.  Schemas, parsers, validators, checkers, versions, routes, claim syntax, evaluators,
budgets, required checks, toolchain choices, and run configuration are \emph{posited}; validator soundness is also
posited whenever it has not been proved.  Parsing and validation invariants, replay implications, arithmetic identities,
certificate composition rules, and any actually established soundness implication are \emph{demonstrated}.  Raw returns,
parse and Boolean outcomes, reports, validation rows, exact values, formatting collisions, floor residues, route
diagnostics, budgets spent, stops, and owned jar walls are \emph{read}.  Exact point ownership, process identity from
equal payloads, claims outside the certificate, cross-toolchain invariance, physical reproducibility, empirical
uncertainty coverage, and deterministic equivalence while nondeterministic inputs remain uncontrolled are
\emph{reserved}.  A report may be useful when validation is partial; it becomes a certificate only for the claims its
checks actually support.
